% ===========================================================================
% alma-micro.tex -- microeconomics notation
% Load with % ===========================================================================
% alma-micro.tex -- microeconomics notation
% Load with % ===========================================================================
% alma-micro.tex -- microeconomics notation
% Load with % ===========================================================================
% alma-micro.tex -- microeconomics notation
% Load with \input{macros/alma-micro}
% ===========================================================================

% --- Sets and spaces -------------------------------------------------------
\providecommand{\R}{\mathbb{R}}          % real numbers
\providecommand{\N}{\mathbb{N}}          % natural numbers
\providecommand{\Q}{\mathbb{Q}}          % rational numbers
\providecommand{\Z}{\mathbb{Z}}          % integers
\newcommand{\Rp}{\mathbb{R}_{+}}         % non-negative reals
\newcommand{\Rpp}{\mathbb{R}_{++}}       % strictly positive reals

% --- Preferences -----------------------------------------------------------
% NOTE ON CORRECTNESS: the symbols below are the standard ones and are easy to
% get backwards. \succsim (weakly preferred to) is the WEAK relation; \succ
% (strictly preferred to) is the STRICT one. The previous Macro.tex had the
% two comments swapped, and named the weak relation \pref and the strict one
% \weakpref -- both misleading. The names here say what they mean.
\newcommand{\weakpref}{\succsim}         % x \weakpref y : x at least as good as y
\newcommand{\strictpref}{\succ}          % x \strictpref y : x strictly better than y
\newcommand{\indiff}{\sim}               % x \indiff y     : indifferent
\newcommand{\weakprefby}{\precsim}       % converse of \weakpref
\newcommand{\strictprefby}{\prec}        % converse of \strictpref
\newcommand{\UCS}{\mathrm{UCS}}          % upper contour set
\newcommand{\LCS}{\mathrm{LCS}}          % lower contour set

% --- Utility and demand ----------------------------------------------------
\newcommand{\uof}[1]{u\!\left(#1\right)}        % utility u(x)
\newcommand{\Ubar}{\overline{u}}                % fixed utility level
\newcommand{\Vof}[1]{v\!\left(#1\right)}        % indirect utility v(p,w)
\newcommand{\eof}[2]{e\!\left(#1,#2\right)}     % expenditure function e(p,u)
\newcommand{\xof}[1]{x\!\left(#1\right)}        % demand x(.)
\newcommand{\Xof}[2]{x\!\left(#1,#2\right)}     % Walrasian demand x(p,w)
\newcommand{\hof}[2]{h\!\left(#1,#2\right)}     % Hicksian demand h(p,u)
\newcommand{\MRS}{\mathrm{MRS}}                 % marginal rate of substitution
\newcommand{\MRT}{\mathrm{MRT}}                 % marginal rate of transformation

% --- Operators -------------------------------------------------------------
\newcommand{\st}{\;\text{s.t.}\;}               % subject to
\newcommand{\limn}{\lim_{n\to\infty}}           % limit as n -> infinity
\newcommand{\argmax}{\operatorname*{arg\,max}}
\newcommand{\argmin}{\operatorname*{arg\,min}}

% --- Formatting helpers ----------------------------------------------------
\newcommand{\abs}[1]{\left|#1\right|}
\newcommand{\norm}[1]{\left\lVert#1\right\rVert}
\newcommand{\set}[1]{\left\{#1\right\}}
\newcommand{\paren}[1]{\left(#1\right)}
\newcommand{\brak}[1]{\left[#1\right]}
\newcommand{\bi}[1]{\boldsymbol{#1}}            % bold vector

% --- Calculus --------------------------------------------------------------
\newcommand{\pder}[2]{\frac{\partial #1}{\partial #2}}
\newcommand{\der}[2]{\frac{d #1}{d #2}}
\newcommand{\lagr}{\mathcal{L}}                 % Lagrangian

% --- Misc ------------------------------------------------------------------
\newcommand{\defeq}{\vcentcolon=}
\newcommand{\eqdef}{=\vcentcolon}

% ===========================================================================

% --- Sets and spaces -------------------------------------------------------
\providecommand{\R}{\mathbb{R}}          % real numbers
\providecommand{\N}{\mathbb{N}}          % natural numbers
\providecommand{\Q}{\mathbb{Q}}          % rational numbers
\providecommand{\Z}{\mathbb{Z}}          % integers
\newcommand{\Rp}{\mathbb{R}_{+}}         % non-negative reals
\newcommand{\Rpp}{\mathbb{R}_{++}}       % strictly positive reals

% --- Preferences -----------------------------------------------------------
% NOTE ON CORRECTNESS: the symbols below are the standard ones and are easy to
% get backwards. \succsim (weakly preferred to) is the WEAK relation; \succ
% (strictly preferred to) is the STRICT one. The previous Macro.tex had the
% two comments swapped, and named the weak relation \pref and the strict one
% \weakpref -- both misleading. The names here say what they mean.
\newcommand{\weakpref}{\succsim}         % x \weakpref y : x at least as good as y
\newcommand{\strictpref}{\succ}          % x \strictpref y : x strictly better than y
\newcommand{\indiff}{\sim}               % x \indiff y     : indifferent
\newcommand{\weakprefby}{\precsim}       % converse of \weakpref
\newcommand{\strictprefby}{\prec}        % converse of \strictpref
\newcommand{\UCS}{\mathrm{UCS}}          % upper contour set
\newcommand{\LCS}{\mathrm{LCS}}          % lower contour set

% --- Utility and demand ----------------------------------------------------
\newcommand{\uof}[1]{u\!\left(#1\right)}        % utility u(x)
\newcommand{\Ubar}{\overline{u}}                % fixed utility level
\newcommand{\Vof}[1]{v\!\left(#1\right)}        % indirect utility v(p,w)
\newcommand{\eof}[2]{e\!\left(#1,#2\right)}     % expenditure function e(p,u)
\newcommand{\xof}[1]{x\!\left(#1\right)}        % demand x(.)
\newcommand{\Xof}[2]{x\!\left(#1,#2\right)}     % Walrasian demand x(p,w)
\newcommand{\hof}[2]{h\!\left(#1,#2\right)}     % Hicksian demand h(p,u)
\newcommand{\MRS}{\mathrm{MRS}}                 % marginal rate of substitution
\newcommand{\MRT}{\mathrm{MRT}}                 % marginal rate of transformation

% --- Operators -------------------------------------------------------------
\newcommand{\st}{\;\text{s.t.}\;}               % subject to
\newcommand{\limn}{\lim_{n\to\infty}}           % limit as n -> infinity
\newcommand{\argmax}{\operatorname*{arg\,max}}
\newcommand{\argmin}{\operatorname*{arg\,min}}

% --- Formatting helpers ----------------------------------------------------
\newcommand{\abs}[1]{\left|#1\right|}
\newcommand{\norm}[1]{\left\lVert#1\right\rVert}
\newcommand{\set}[1]{\left\{#1\right\}}
\newcommand{\paren}[1]{\left(#1\right)}
\newcommand{\brak}[1]{\left[#1\right]}
\newcommand{\bi}[1]{\boldsymbol{#1}}            % bold vector

% --- Calculus --------------------------------------------------------------
\newcommand{\pder}[2]{\frac{\partial #1}{\partial #2}}
\newcommand{\der}[2]{\frac{d #1}{d #2}}
\newcommand{\lagr}{\mathcal{L}}                 % Lagrangian

% --- Misc ------------------------------------------------------------------
\newcommand{\defeq}{\vcentcolon=}
\newcommand{\eqdef}{=\vcentcolon}

% ===========================================================================

% --- Sets and spaces -------------------------------------------------------
\providecommand{\R}{\mathbb{R}}          % real numbers
\providecommand{\N}{\mathbb{N}}          % natural numbers
\providecommand{\Q}{\mathbb{Q}}          % rational numbers
\providecommand{\Z}{\mathbb{Z}}          % integers
\newcommand{\Rp}{\mathbb{R}_{+}}         % non-negative reals
\newcommand{\Rpp}{\mathbb{R}_{++}}       % strictly positive reals

% --- Preferences -----------------------------------------------------------
% NOTE ON CORRECTNESS: the symbols below are the standard ones and are easy to
% get backwards. \succsim (weakly preferred to) is the WEAK relation; \succ
% (strictly preferred to) is the STRICT one. The previous Macro.tex had the
% two comments swapped, and named the weak relation \pref and the strict one
% \weakpref -- both misleading. The names here say what they mean.
\newcommand{\weakpref}{\succsim}         % x \weakpref y : x at least as good as y
\newcommand{\strictpref}{\succ}          % x \strictpref y : x strictly better than y
\newcommand{\indiff}{\sim}               % x \indiff y     : indifferent
\newcommand{\weakprefby}{\precsim}       % converse of \weakpref
\newcommand{\strictprefby}{\prec}        % converse of \strictpref
\newcommand{\UCS}{\mathrm{UCS}}          % upper contour set
\newcommand{\LCS}{\mathrm{LCS}}          % lower contour set

% --- Utility and demand ----------------------------------------------------
\newcommand{\uof}[1]{u\!\left(#1\right)}        % utility u(x)
\newcommand{\Ubar}{\overline{u}}                % fixed utility level
\newcommand{\Vof}[1]{v\!\left(#1\right)}        % indirect utility v(p,w)
\newcommand{\eof}[2]{e\!\left(#1,#2\right)}     % expenditure function e(p,u)
\newcommand{\xof}[1]{x\!\left(#1\right)}        % demand x(.)
\newcommand{\Xof}[2]{x\!\left(#1,#2\right)}     % Walrasian demand x(p,w)
\newcommand{\hof}[2]{h\!\left(#1,#2\right)}     % Hicksian demand h(p,u)
\newcommand{\MRS}{\mathrm{MRS}}                 % marginal rate of substitution
\newcommand{\MRT}{\mathrm{MRT}}                 % marginal rate of transformation

% --- Operators -------------------------------------------------------------
\newcommand{\st}{\;\text{s.t.}\;}               % subject to
\newcommand{\limn}{\lim_{n\to\infty}}           % limit as n -> infinity
\newcommand{\argmax}{\operatorname*{arg\,max}}
\newcommand{\argmin}{\operatorname*{arg\,min}}

% --- Formatting helpers ----------------------------------------------------
\newcommand{\abs}[1]{\left|#1\right|}
\newcommand{\norm}[1]{\left\lVert#1\right\rVert}
\newcommand{\set}[1]{\left\{#1\right\}}
\newcommand{\paren}[1]{\left(#1\right)}
\newcommand{\brak}[1]{\left[#1\right]}
\newcommand{\bi}[1]{\boldsymbol{#1}}            % bold vector

% --- Calculus --------------------------------------------------------------
\newcommand{\pder}[2]{\frac{\partial #1}{\partial #2}}
\newcommand{\der}[2]{\frac{d #1}{d #2}}
\newcommand{\lagr}{\mathcal{L}}                 % Lagrangian

% --- Misc ------------------------------------------------------------------
\newcommand{\defeq}{\vcentcolon=}
\newcommand{\eqdef}{=\vcentcolon}

% ===========================================================================

% --- Sets and spaces -------------------------------------------------------
\providecommand{\R}{\mathbb{R}}          % real numbers
\providecommand{\N}{\mathbb{N}}          % natural numbers
\providecommand{\Q}{\mathbb{Q}}          % rational numbers
\providecommand{\Z}{\mathbb{Z}}          % integers
\newcommand{\Rp}{\mathbb{R}_{+}}         % non-negative reals
\newcommand{\Rpp}{\mathbb{R}_{++}}       % strictly positive reals

% --- Preferences -----------------------------------------------------------
% NOTE ON CORRECTNESS: the symbols below are the standard ones and are easy to
% get backwards. \succsim (weakly preferred to) is the WEAK relation; \succ
% (strictly preferred to) is the STRICT one. The previous Macro.tex had the
% two comments swapped, and named the weak relation \pref and the strict one
% \weakpref -- both misleading. The names here say what they mean.
\newcommand{\weakpref}{\succsim}         % x \weakpref y : x at least as good as y
\newcommand{\strictpref}{\succ}          % x \strictpref y : x strictly better than y
\newcommand{\indiff}{\sim}               % x \indiff y     : indifferent
\newcommand{\weakprefby}{\precsim}       % converse of \weakpref
\newcommand{\strictprefby}{\prec}        % converse of \strictpref
\newcommand{\UCS}{\mathrm{UCS}}          % upper contour set
\newcommand{\LCS}{\mathrm{LCS}}          % lower contour set

% --- Utility and demand ----------------------------------------------------
\newcommand{\uof}[1]{u\!\left(#1\right)}        % utility u(x)
\newcommand{\Ubar}{\overline{u}}                % fixed utility level
\newcommand{\Vof}[1]{v\!\left(#1\right)}        % indirect utility v(p,w)
\newcommand{\eof}[2]{e\!\left(#1,#2\right)}     % expenditure function e(p,u)
\newcommand{\xof}[1]{x\!\left(#1\right)}        % demand x(.)
\newcommand{\Xof}[2]{x\!\left(#1,#2\right)}     % Walrasian demand x(p,w)
\newcommand{\hof}[2]{h\!\left(#1,#2\right)}     % Hicksian demand h(p,u)
\newcommand{\MRS}{\mathrm{MRS}}                 % marginal rate of substitution
\newcommand{\MRT}{\mathrm{MRT}}                 % marginal rate of transformation

% --- Operators -------------------------------------------------------------
\newcommand{\st}{\;\text{s.t.}\;}               % subject to
\newcommand{\limn}{\lim_{n\to\infty}}           % limit as n -> infinity
\newcommand{\argmax}{\operatorname*{arg\,max}}
\newcommand{\argmin}{\operatorname*{arg\,min}}

% --- Formatting helpers ----------------------------------------------------
\newcommand{\abs}[1]{\left|#1\right|}
\newcommand{\norm}[1]{\left\lVert#1\right\rVert}
\newcommand{\set}[1]{\left\{#1\right\}}
\newcommand{\paren}[1]{\left(#1\right)}
\newcommand{\brak}[1]{\left[#1\right]}
\newcommand{\bi}[1]{\boldsymbol{#1}}            % bold vector

% --- Calculus --------------------------------------------------------------
\newcommand{\pder}[2]{\frac{\partial #1}{\partial #2}}
\newcommand{\der}[2]{\frac{d #1}{d #2}}
\newcommand{\lagr}{\mathcal{L}}                 % Lagrangian

% --- Misc ------------------------------------------------------------------
\newcommand{\defeq}{\vcentcolon=}
\newcommand{\eqdef}{=\vcentcolon}
